% =============================================================
% commands.tex — Commandes et environnements personnalisés
% =============================================================

% --- Deux sorties depuis les mêmes sources ---
% `main.tex` compile par défaut la version de remise. Le petit point
% d'entrée `main_brouillon.tex` définit \menalbrouillon avant de charger
% `main.tex`, ce qui rend alors visibles les informations encore ouvertes.
\newif\ifbrouillon
\ifdefined\menalbrouillon
  \brouillontrue
\else
  \brouillonfalse
\fi

% --- Colonne de tableau en drapeau (non justifiée), utilisée pour toutes
%     les colonnes p{...} : évite les "Underfull \hbox" massifs dans les
%     colonnes étroites et les ruptures de page problématiques en fin de
%     longtable ("Infinite glue shrinkage") causées par la justification
%     forcée sur peu de mots par ligne ---
\newcolumntype{P}[1]{>{\raggedright\arraybackslash}p{#1}}

% --- Marqueur visible pour tout contenu manquant ---
% Usage : \todo{ce qui manque et pourquoi}
\newcommand{\todo}[1]{%
  \ifbrouillon
    \textcolor{red!70!black}{\textbf{[À COMPLÉTER --- #1]}}%
  \fi
}

% --- Encadré pour les sections bloquées par un document source
%     manquant (cf. 00_SOMMAIRE_PFE_MENAL.md §5) ---
\newtcolorbox{attentedoc}[1][]{
  colback=red!5, colframe=red!60!black,
  colbacktitle=red!60!black, coltitle=white,
  fonttitle=\bfseries,
  title={En attente de document source},
  breakable, #1
}

% --- Encadré "principe clé" / point à retenir ---
\newtcolorbox{keybox}[1][]{
  colback=blue!5, colframe=blue!40!black,
  colbacktitle=blue!40!black, coltitle=white,
  fonttitle=\bfseries,
  breakable, #1
}

% --- Encadré d'alerte / point de vigilance ---
\newtcolorbox{alertbox}[1][]{
  colback=orange!8, colframe=orange!70!black,
  colbacktitle=orange!70!black, coltitle=white,
  fonttitle=\bfseries,
  title={Point de vigilance},
  breakable, #1
}

% --- Encadré "extrait de configuration" (courts extraits encadrés,
%     jamais de fichier de code complet dans le corps du rapport) ---
\newtcolorbox{extraitconf}[1][]{
  colback=gray!5, colframe=gray!50!black,
  fonttitle=\bfseries\ttfamily\small, coltitle=black,
  fontupper=\ttfamily\scriptsize,
  breakable, #1
}

% --- Renvoi rapide vers le sommaire de travail (usage interne,
%     à retirer avant impression finale) ---
\newcommand{\refsommaire}[1]{\textit{(cf. sommaire de travail, #1)}}

% --- Emplacement d'une capture d'écran ou d'un schéma pas encore
%     produit, à l'intérieur d'un environnement figure/table/graphique.
%     Usage : \figureaproduire{ce qui doit apparaître, et ce qui doit
%     être masqué -- cf. PLAN_DES_FIGURES_ET_CAPTURES.md §1} ---
\newcommand{\figureaproduire}[1]{%
  \ifbrouillon
    \fbox{\parbox{0.92\linewidth}{\centering\vspace{1.4cm}\todo{#1}\vspace{1.4cm}}}%
  \fi
}

% --- Citation dont la clé .bib n'existe pas encore (bibliographie en
%     cours de consolidation, cf. bibliography/references.bib). Usage :
%     \citetodo{Auteurs, Titre, année} --- à remplacer par \cite{clé}
%     une fois la référence ajoutée au fichier .bib. ---
\newcommand{\citetodo}[1]{%
  \ifbrouillon
    \textcolor{red!70!black}{[réf. à ajouter --- #1]}%
  \fi
}

% --- Date de référence unique du rapport ---
% Tous les statuts sont exprimés « à la date de \dateref ».
\newcommand{\dateref}{25/08/2026}

% --- Marqueur de contenu à compléter par l'auteur ---
% Visible en mode brouillon, invisible en mode final.
\newcommand{\acompleter}[1]{%
  \ifbrouillon
    \textcolor{red!70!black}{\textbf{[À COMPLÉTER --- #1]}}%
  \else
    % En mode final, le placeholder est invisible
  \fi
}

% --- Marqueur de contenu en attente d'une décision ---
\newcommand{\attente}[1]{%
  \ifbrouillon
    \textcolor{orange!80!black}{\textit{[En attente --- #1]}}%
  \fi
}

% --- Capture d'écran : emplacement réservé (norme ESPRIT §14) ---
% Usage : \capture{ID}{légende}{consigne de masquage}
%         \capture[7cm]{ID}{légende}{consigne de masquage}
% Si le fichier figures/<ID>.png existe, il est inséré ; sinon un cadre
% de hauteur réservée tient la place, avec l'identifiant en petit gris
% en version de remise et la consigne de masquage en version de travail.
\newcommand{\captureheightdefault}{6cm}
\newcommand{\capture}[4][\captureheightdefault]{%
  \begin{figure}[htbp]
    \centering
    \IfFileExists{figures/#2.png}{%
      \setlength{\fboxrule}{0.4pt}\setlength{\fboxsep}{0pt}%
      \fcolorbox{black!30}{white}{\includegraphics[width=0.80\linewidth]{figures/#2.png}}%
    }{%
      \setlength{\fboxrule}{0.4pt}\setlength{\fboxsep}{0pt}%
      \fcolorbox{black!30}{white}{%
        \parbox[c][#1][c]{0.88\linewidth}{%
          \centering
          \ifbrouillon
            \footnotesize\color{red!70!black}%
            \textbf{[CAPTURE #2 --- À PRODUIRE]}\par\vspace{2pt}%
            \begin{minipage}{0.90\linewidth}\centering\itshape #4\end{minipage}%
          \else
            \scriptsize\color{black!45}\textsf{Capture #2}%
          \fi
        }%
      }%
    }%
    \caption{#3}
    \label{fig:capture-#2}
  \end{figure}%
}

% --- Capture double a/b : deux preuves d'un même test ---
% \capturedouble{IDa}{IDb}{légende globale}{sous-lég. a}{sous-lég. b}{consigne}
\newcommand{\capturedouble}[6]{%
  \begin{figure}[htbp]
    \centering
    \begin{subfigure}[b]{0.48\linewidth}\centering
      \IfFileExists{figures/#1.png}%
        {\fbox{\includegraphics[width=0.94\linewidth]{figures/#1.png}}}%
        {\fcolorbox{black!30}{white}{\parbox[c][3.4cm][c]{0.92\linewidth}{\centering
           \ifbrouillon\footnotesize\color{red!70!black}\textbf{[CAPTURE #1]}%
           \else\scriptsize\color{black!45}\textsf{Capture #1}\fi}}}%
      \caption{#4}\label{fig:capture-#1}
    \end{subfigure}\hfill
    \begin{subfigure}[b]{0.48\linewidth}\centering
      \IfFileExists{figures/#2.png}%
        {\fbox{\includegraphics[width=0.94\linewidth]{figures/#2.png}}}%
        {\fcolorbox{black!30}{white}{\parbox[c][3.4cm][c]{0.92\linewidth}{\centering
           \ifbrouillon\footnotesize\color{red!70!black}\textbf{[CAPTURE #2]}%
           \else\scriptsize\color{black!45}\textsf{Capture #2}\fi}}}%
      \caption{#5}\label{fig:capture-#2}
    \end{subfigure}
    \caption{#3}
    \label{fig:capture-#1#2}
    \ifbrouillon\par\vspace{2pt}{\footnotesize\color{red!70!black}\itshape Masquage : #6}\fi
  \end{figure}%
}

% --- Horodatage : neutralise l'espace fine française avant les deux-points ---
% Usage : \hms{16:41:05}
\DeclareRobustCommand{\hms}[1]{{\NoAutoSpaceBeforeFDP #1}}

% --- Pictogrammes de couverture (cercle plein / mi-plein / vide) ---
%     Employes au chapitre 1 et dans l'annexe A. Dessines en TikZ plutot
%     qu'en glyphes Unicode, dont la couverture est incertaine en Times.
\DeclareRobustCommand{\pictoPleinAudit}{\tikz[baseline=-0.5ex]{\fill[black] (0,0) circle (0.07cm);}}
\DeclareRobustCommand{\pictoPartielAudit}{\tikz[baseline=-0.5ex]{\fill[black] (0,0) -- ++(90:0.07cm) arc (90:270:0.07cm) -- cycle; \draw[black] (0,0) circle (0.07cm);}}
\DeclareRobustCommand{\pictoVideAudit}{\tikz[baseline=-0.5ex]{\draw[black] (0,0) circle (0.07cm);}}


% --- Encadré « Résultat mesuré » ---
% Réservé aux valeurs RÉELLEMENT relevées et datées. Une dizaine dans tout
% le mémoire, pas davantage : c'est ce qui les rend repérables au feuilletage.
% Usage : \begin{mesurebox}{titre court} ... \end{mesurebox}
\newtcolorbox{mesurebox}[1]{
  colback=green!4, colframe=green!35!black,
  colbacktitle=green!35!black, coltitle=white,
  fonttitle=\bfseries,
  title={Résultat mesuré --- #1},
  breakable
}


% --- Symbole euro ---
% XeLaTeX avec une police système : le glyphe est disponible directement.
% Défini ici plutôt que chargé par un paquet, pour ne pas ajouter de dépendance.
\providecommand{\euro}{€}

% --- Convention graphique unique (styles TikZ partages par toutes les figures) ---
% =============================================================
% CONVENTION GRAPHIQUE UNIQUE DU MÉMOIRE
%
% Un aspect de trait = une signification, stable dans tout le document.
% Une couleur = une signification. L'accent « refus » ne porte jamais seul
% le sens : un refus est toujours aussi un trait épais ET une croix, de
% sorte que la figure reste lisible à l'impression en noir et blanc.
%
% Bibliothèques requises : arrows.meta, positioning, shapes.geometric,
% calc, fit, backgrounds — toutes déjà chargées par config/packages.tex.
% Aucun paquet nouveau.
% =============================================================

% --- 1. PALETTE : cinq teintes, une signification chacune -----------------
\definecolor{menalencre}{gray}{0.15}   % traits et libellés porteurs
\definecolor{menaltrait}{gray}{0.45}   % traits secondaires, annotations
\definecolor{menalplanA}{gray}{0.970}  % bande — plan de données
\definecolor{menalplanB}{gray}{0.900}  % bande — plan de contrôle
\definecolor{menalplanC}{gray}{0.830}  % bande — plan d'observation
\definecolor{menalrefus}{rgb}{0.55,0.10,0.10} % refus / faiblesse (accent unique)

% --- 2. STYLES DE TRAIT ---------------------------------------------------
\tikzset{
  fluxdonnee/.style   = {-{Latex[length=1.7mm]}, line width=0.6pt,
                         draw=menalencre},
  fluxreponse/.style  = {-{Latex[length=1.7mm,open]}, line width=0.6pt,
                         draw=menaltrait},
  fluxevenement/.style= {-{Latex[length=1.7mm]}, line width=0.6pt,
                         draw=menalencre, densely dashed},
  fluxcontrole/.style = {-{Latex[length=1.7mm]}, line width=0.6pt,
                         draw=menalencre, densely dotted},
  fluxrefuse/.style   = {-{Latex[length=1.9mm]}, line width=1.4pt,
                         draw=menalrefus, line cap=round},
  lienstructure/.style= {line width=0.5pt, draw=menaltrait},
  % Pas de style « dépendance externe » : essayé, puis écarté --- un trait plein
  % à pointe creuse est indistinguable de fluxdonnee en niveaux de gris.
  % L'extériorité est portée par la forme du nœud (nexterne).
  fluxsynthese/.style = {-{Latex[length=2.4mm,width=1.8mm]}, line width=1.6pt,
                         draw=menalencre, line join=round},
  seuil/.style        = {line width=1.0pt, draw=menalencre, dashed},
  frontiereplan/.style= {line width=1.6pt, draw=menalencre},
  % Croix de refus : \draw[fluxrefuse] (a) -- node[croixrefus]{} (b);
  % Le fond blanc masque le trait sous la croix.
  croixrefus/.style   = {pos=0.5, inner sep=0pt, minimum size=2.4mm,
                         draw=none, fill=white,
                         path picture={%
                           \draw[line width=1.1pt, menalrefus, line cap=round]
                             (path picture bounding box.south west)
                               -- (path picture bounding box.north east)
                             (path picture bounding box.north west)
                               -- (path picture bounding box.south east);}},
  etiq/.style         = {midway, fill=white, inner sep=1.2pt,
                         font=\footnotesize, text=menalencre},
  etiqsec/.style      = {midway, fill=white, inner sep=1.2pt,
                         font=\scriptsize, text=menaltrait},
}

% --- 3. FORMES DE NŒUD ----------------------------------------------------
\tikzset{
  nsysteme/.style  = {draw=menalencre, line width=0.6pt, rounded corners=2pt,
                      fill=white, align=center, font=\footnotesize,
                      inner sep=3.5pt, minimum height=0.80cm},
  nexterne/.style  = {nsysteme, densely dashed},
  nstock/.style    = {nsysteme, double, double distance=0.8pt,
                      rounded corners=0pt},
  nacteur/.style   = {nsysteme, rounded corners=6pt, fill=menalplanA},
  nporte/.style    = {draw=menalencre, line width=1.1pt, diamond,
                      aspect=2.2, fill=menalplanB, align=center,
                      font=\footnotesize, inner sep=2pt},
  ncontrole/.style = {nsysteme, line width=1.2pt},
  nabsent/.style   = {nsysteme, densely dashed, draw=menaltrait,
                      text=menaltrait},
  bandedonnee/.style      = {fill=menalplanA, rounded corners=3pt,
                             inner sep=5pt, draw=none},
  bandecontrole/.style    = {fill=menalplanB, rounded corners=3pt,
                             inner sep=5pt, draw=none},
  bandeobservation/.style = {fill=menalplanC, rounded corners=3pt,
                             inner sep=5pt, draw=none},
  titrebande/.style= {font=\footnotesize\itshape, text=menalencre,
                      inner sep=2pt},
}

% --- 4. LÉGENDE INTERNE : obligatoire dès deux aspects de trait -----------
% Déclarées ROBUSTES, pour la même raison que les pictogrammes de couverture :
% une commande TikZ fragile casse dès qu'elle remonte dans une légende ou
% dans la liste des figures.
% Les deux macros n'imposent aucune taille : elles héritent de celle du contexte.
% Une taille forcée composait « refus » plus gros que ses voisins dans la ligne de
% notation d'une figure, qui est en \scriptsize.
\DeclareRobustCommand{\legendetrait}[2]{%
  \tikz[baseline=-0.55ex]{\draw[#1] (0,0) -- (0.7,0);}~{#2}%
}
\DeclareRobustCommand{\legendecroix}[1]{%
  \tikz[baseline=-0.55ex]{\draw[fluxrefuse] (0,0) -- node[croixrefus]{} (0.7,0);}%
  ~{#1}%
}
\tikzset{
  legendecadre/.style = {draw=menaltrait, line width=0.4pt, rounded corners=2pt,
                         fill=white, inner sep=4pt, align=left,
                         font=\footnotesize},
}

% --- 5. ANNOTATION SECONDAIRE --------------------------------------------
% Nom de produit ou compteur cité dans un nœud, sans valeur démonstrative.
% Défini ici plutôt que dans un chapitre : plusieurs chapitres l'emploient.
\providecommand{\gcprod}[1]{{\scriptsize\color{menaltrait}#1}}

